% Method
\section{EGC-TTRL}
EGC-TTRL instantiates one principle: \emph{uncertain deployment-time
evidence should not unconditionally change the future policy}. Two gates
implement it at two scales: a \textbf{local gate} decides which action
tokens may produce gradient; a \textbf{global gate} decides whether the
resulting candidate adapter may reach future tasks.

\subsection{Local gate: paired-branch signed credit}
At a selected decision state $s_t$, we restore the sealed snapshot and
sample $G=4$ grammar-valid alternative actions from the parent policy,
each continued under $R=2$ common-random-number seeds with the same
continuation protocol. Executing each branch in the environment yields
an evidence utility matrix $U\in[0,1]^{G\times R}$; the signed credit of
action $i$ is
\begin{equation}
\hat c_i=\bar U_i-\tfrac{1}{G}\sum_j \bar U_j,
\qquad
z_{i,r}=U_{i,r}-\tfrac{1}{G}\sum_j U_{j,r},
\end{equation}
with a reliability interval $[\hat c_i-b_i,\hat c_i+b_i]$ where
$b_i=t_{R-1,0.90}\sqrt{\hat v_i/R}$ and $\hat v_i$ is the across-seed
variance of $z$. The local gate passes action $i$ iff its interval is
disjoint from $[-\eta_{\mathrm{credit}},+\eta_{\mathrm{credit}}]$;
credit is clipped and group-normalized. An alternative local gate uses
\textbf{evidence conflict}: when hard evidence contradicts the soft
verifier (e.g., a poisoned receipt vs.\ a state invariant), the whole
group abstains and a drift monitor may halt adaptation.

\subsection{Action-token objective}
Each passed action yields one immutable update row over its own
generated action-token span (never observation/prefix tokens), with the
behavior log-probability from the generation service. The objective is a
clipped policy ratio with KL-to-parent regularization (the baseline
naive variant uses the same objective with group-relative terminal
utilities and full-completion masks; EGC differs only in the credit
signal).

\subsection{Global gate: SafeCommit}
Each update produces a non-servable candidate adapter $\phi'_k$. A
shadow evaluation compares candidate vs.\ parent on fresh sentinel gain
tasks and anchor capability tasks, with per-candidate error budgets
$\alpha_k=6\alpha_{\mathrm{total}}/(\pi^2 k^2)$. The commit decision
uses an empirical-Bernstein e-process confidence sequence (frozen by a
pre-registered coverage simulator over 162 configurations, where the
fixed-n Hoeffding alternative passed no operating point):
\begin{equation}
\mathrm{LCB}_{\mathrm{gain}}\ge \varepsilon_{\mathrm{gain}}
\;\wedge\;
\mathrm{UCB}_{\mathrm{harm}}\le \varepsilon_{\mathrm{harm}}
\;\wedge\;
\mathrm{Guard}=\mathrm{ALLOW}.
\end{equation}
Otherwise the candidate is rolled back atomically. Sentinel items are
used at most once; coverage claims are restricted to the proxy-mean
decision error.

\subsection{Statistics}
Primary contrasts use paired hierarchical bootstrap over
seed/domain-stream and task families, with pre-registered tests at
$p<0.01$ plus effect sizes; the number of seeds is frozen before test
data is seen.
